\chapter{Results and Discussion}
\label{ch:results}

\section{Path Following Performance}

% TODO: training curves (episode return vs. timesteps), mean/std CTE per model

\section{Reverse Maneuvering Success}

% TODO: jackknife rate per model, safe articulation zone analysis, trajectory plots

\section{Perception Robustness}

% TODO: CTE degradation under sensor noise, F1-score vs. noise level curves

\section{Benchmark Comparison}
% Controller benchmark — best TD3 configuration vs tuned classical controllers.
% All methods run on identical pre-generated scenario sets (fixed seeds).
% Classical controllers are tuned via differential evolution before comparison.

\begin{table}[H]
\centering
\caption{Forward lane-following: controller benchmark (no obstacles, 50~eval episodes,
         held-out seeds).}
\label{tab:benchmark_forward}
\begin{tabular}{lcccccc}
\toprule
Controller
  & \makecell{CTE\\Trailer\\mean (m)}
  & \makecell{CTE\\Trailer\\RMS (m)}
  & \makecell{Max $|\gamma|$\\($^\circ$)}
  & \makecell{Completion\\(\%)}
  & \makecell{Jackknife\\(\%)}
  & \makecell{Latency\\(ms)} \\
\midrule
TD3 (best config)      & --- & --- & --- & --- & --- & --- \\
Forward pure pursuit   & --- & --- & --- & --- & --- & $<$1 \\
PID (tuned)            & --- & --- & --- & --- & --- & $<$1 \\
MPC (tuned)            & --- & --- & --- & --- & --- & --- \\
\bottomrule
\end{tabular}
\end{table}

\begin{table}[H]
\centering
\caption{Reverse lane-following: controller benchmark (no obstacles, 50~eval episodes,
         held-out seeds).}
\label{tab:benchmark_reverse}
\begin{tabular}{lcccccc}
\toprule
Controller
  & \makecell{CTE\\Trailer\\mean (m)}
  & \makecell{CTE\\Trailer\\RMS (m)}
  & \makecell{Max $|\gamma|$\\($^\circ$)}
  & \makecell{Completion\\(\%)}
  & \makecell{Jackknife\\(\%)}
  & \makecell{Latency\\(ms)} \\
\midrule
TD3 (best config)      & --- & --- & --- & --- & --- & --- \\
Reverse pure pursuit   & --- & --- & --- & --- & --- & $<$1 \\
PID reverse (tuned)    & --- & --- & --- & --- & --- & $<$1 \\
MPC reverse (tuned)    & --- & --- & --- & --- & --- & --- \\
\bottomrule
\end{tabular}
\end{table}


The results must be understood in terms of multi-objective performance: in rigid-body systems, success is synonymous with minimising cross-track error. In articulated systems, hitch stability takes precedence over exact path following. An agent that maintains $|\gamma| < 5^\circ$ throughout a reverse task is more successful than one that achieves lower CTE but triggers jackknife events.

% TODO: statistical significance (confidence intervals), ablation over reward weight choices
\begin{table}[H]
\centering
\caption{Recommended experiment matrix for the thesis}
\label{tab:experiment_matrix}
\small
\begin{tabular}{>{\raggedright\arraybackslash}p{3.6cm} >{\raggedright\arraybackslash}p{2.5cm} >{\raggedright\arraybackslash}p{2.7cm} >{\raggedright\arraybackslash}p{4.6cm}}
\toprule
Experiment & Task & Observation & Success Metric \\
\midrule
Observation ablation (forward) & Path following & state, LiDAR, BEV variants & Fixed-speed trailer mean abs.\ CTE, max $|\gamma|$, completion \\
Observation ablation (reverse) & Reverse tracking & state, LiDAR, BEV variants & Fixed-speed trailer mean abs.\ CTE, max $|\gamma|$, jackknife rate \\
Reward ablation (forward) & Path following & best observation only & Fixed-speed mean return, trailer CTE, stability metrics \\
Reward ablation (reverse) & Reverse tracking & best observation only & Fixed-speed completion, jackknife rate, trailer CTE \\
Controller benchmark (forward) & Path following & best RL obs.\ vs.\ classical & Fixed-speed TD3 vs.\ pure pursuit vs.\ PID vs.\ MPC \\
Controller benchmark (reverse) & Reverse tracking & best RL obs.\ vs.\ classical & Fixed-speed TD3 vs.\ reverse pure pursuit vs.\ reverse PID vs.\ reverse MPC \\
Obstacle RL extension & End-to-end navigation & LiDAR or BEV & Variable-speed collision rate, completion, min clearance \\
Obstacle classical comparison & Planner + controller & matched local path + controller & Only if classical baselines receive equivalent planning information \\
\bottomrule
\end{tabular}
\end{table}

